\documentclass{article}
\usepackage{draftwatermark}
\usepackage{hyperref}
\usepackage[type={CC}, modifier={zero}, version={1.0}]{doclicense}
\usepackage{amsmath}
%\usepackage[bb = boondox]{mathalpha}

\title{Tensor Products, Exterior Products, and Geometric Algebras Aren't Trivial}
\author{Jacob Bourbaki}


\begin{document}
\SetWatermarkText{DRAFT}
\SetWatermarkColor{gray!75}
\SetWatermarkScale{5}
\maketitle
\doclicenseThis
\begin{abstract}
Tensor, exterior, and geometric algebras can be nicely defined in terms
of free algebras with relations.  However, this begs the question of whether, once
all of the indentifications have been made, what elements remain distinct from zero.
This note is a comment on how to check that the identifications don't identify more
that one wanted.  There is also a comment on understanding the determinant.  There is
no claim to originality, although no references seem to be immediately available.
\end{abstract}

\section{Introduction}
It is well known that tensort product, exterior products, and geometric products
can be defined in terms of generators and relations.  However, such definitions
leave open the \emph{a priori} possibility that the algebra so defined is degenerate
due to relationships which are not obvious.  For example, if one considers
the groups $\mathbf{Z}/2\mathbf{Z}$ and $\mathbf{Z}/3\mathbf{Z}$ as 
$\mathbf{Z}$--modules, then
\[
  \mathbf{Z}/2\mathbf{Z}  \otimes_{\mathbf{Z}} \mathbf{Z}/3\mathbf{Z}  = 0,
\]
that is, all of the elements are identically equal to zero, as all elements
of $\mathbf{Z}/2\mathbf{Z}$ are divisible by $3=1$~(mod~2), but when the 3~moves
accross the $\otimes_{\mathbf{Z}}$, it becomes zero mod 3.  Thus the only 
element of the tensor product is zero.  Fortunately, this doesn't have for 
fields like the reals, $\mathbf{R}$.  The purpose of this note is to make this
obvious.

\section{Tensor Algebra}
For a vector space $V$ over a field $\mathbf{F}$, let $W=V^*$ be the dual space.
Now, for each integer $n>0$, let $X^n$ the vector space whose elements are
lists of $n$ vectors, so $\left[w_1, w_2, \ldots w_n\right]$, $w_i \in W$, with addition
and multiplication defined component--wise.  Let $X^0 = \mathbf{F}$, and let
\begin{equation}
  \mathbf{X}  = \bigoplus_{i=0}^\infty X^i.
\end{equation}
Then $\mathbf{X}$ is a vector space of $\mathbf{F}$, and we can consider the alegbra of
all $\mathbf{F}$--linear operators from $\mathbf{X}$ to itself, \emph{i.e.}, 
$\text{Hom}(\mathbf{X}, \mathbf{X})$.
Given any $v \in V$, one can associate a linear operator $\hat{v}$ defined by
\begin{align}
  \hat{v}(x) & = 0 & x \in \mathbf{F} \\
  \hat{v}\left(\left[ w_1, w_2, \ldots, w_n\right]\right) 
      &= w_1(v) \left(\left[ w_2, \ldots, w_n\right]\right)  &
           \left[ w_1, w_2, \ldots, w_n\right] \in X^n
\end{align}
Thus each $v\in V$ induces a linear map, and this can be clearly extended to $V \times V$ is
a manner which is probably balanced, and for any basis of $V$, the ordered pairs of
vectors give a basis for $V\otimes V$, \emph{etc}. (Maybe a bit more should be said\ldots.)

\section{Exterior Product}
Using the same definitions for $X^i$ and $\mathbf{X}$ as in the section for tensor products,
define
\begin{align}
  \tilde{v}(x) & = 0 \quad x \in \mathbf{F} \\
  \tilde{v}\left(\left[ w_1, w_2, \ldots, w_n\right]\right) 
      &= w_1(v) \left(\left[ w_2, \ldots, w_n\right]\right)    \nonumber \\
      &- w_2(v) \left(\left[ w_1, w_3, \ldots, w_n\right]\right)   \nonumber \\
      &+ \ldots  \nonumber \\
      &+ (-1)^{n+1} w_n(v) \left(\left[ w_1, w_3, \ldots, w_{n-1}\right]\right)   \nonumber \\
      & \quad \quad \quad \quad \quad     \left[ w_1, w_2, \ldots, w_n\right] \in X^n
\end{align}
Again, this embeds each $v\in V$ in the algebra $\text{Hom}(\mathbf{X}, \mathbf{X})$
in such a way that
\[
    \tilde{v}_1 \tilde{v}_2 + \tilde{v}_2 \tilde{v}_1 = 0.
\]
Thus that exterior algebra is isomophically represented by the subalgebra of 
$\text{Hom}(\mathbf{X}, \mathbf{X})$ generated by the images of the elements of $V$. 

Further, clearly the grading of the algebra is preserved, \emph{etc.}:
\[
  \bigwedge V = \bigwedge_{i=0}^{\infty} V^i
\]


Also note that if multiplying elements of the algebra we don't need to say anything about
factors of $n!$.


Note that for each $w\in W=V^*$, this is well defined:
\begin{align}
  i_w (v_1 \wedge v_2 \wedge \ldots \wedge v_k) \nonumber \\
  \quad = w(v_1) v_2 \wedge \ldots \wedge v_k) \nonumber \\
  \quad \quad - w(v_2) v_1 \wedge v_3 \ldots \wedge v_k) \nonumber \\
  \quad \quad + \ldots  \nonumber \\
  \quad \quad + (-1)^{k+1} w(v_k) v_1 \wedge v_2  \wedge \ldots \wedge v_{k-1}) 
\end{align}


\subsection{Determinants}
If $V=\mathbf{F}^n$, then treating $v_i\in V$ and column vectors and letting 
$\omega_i \in W=V^*$ be the standard dual basis, then the determinant is
\[
\left| \: v_1 \: v_2 \: \ldots \: v_n \right| 
 = \tilde{v}_1 \tilde{v}_2 \ldots \tilde{v}_n \left( \left[ \omega_n, \omega_{n-1}, \ldots,
                                                      \omega_1 \right] \right).
\]

Considering the determinant as a signed volume, note that the determinat is actually fixed
up to a scalar constant by the two conditions:
\begin{enumerate}
\item If $v_i$ is replaced with $\alpha v_i$, $\alpha \in \mathbf{f}$, 
then the determinant is changed by a factor of $\alpha$.
\item If $v_i$ is replaced by $v_i + \alpha v_j$ for $i \ne j$, the value of the 
determinant does not change; this corresponds to the idea that a shear parallel to
one of the faces of a polytope does not change the volume of the polytope.
\end{enumerate}


(Yes, I should say more \ldots.)
\section{Geometric Algebra}
Proceeding to the geometric algebra, let $g(v,v')\in \mathbf{F}$ be a symmetric bilinear form
defined on each pair $v, v' \in V$.  So for each $v \in V$, one has $g(v,\cdot) \in V^*$.
Consider $\bigwedge V$ as a vector space, for each $v \in V$ one can associate linear endomorphism
on $\bigwedge V$ by 
\[
  v^\dagger (v_1 \wedge \ldots \wedge v_k) = \pm i_{g(v,\cdot)} (v_1 \wedge \ldots v_k)
                                          \pm v \wedge v_1 \wedge \ldots v_k
\]

(or something like that... that's all for now).



\end{document}
